\chapter{Conclusão}
\label{cap:end}

O desenvolvimento da JagozApp permitiu construir uma solução web orientada à gestão de processos administrativos e operacionais de um clube desportivo local. Ao longo do projeto foi possível centralizar num único sistema funcionalidades que, tradicionalmente, tendem a estar dispersas por diferentes ferramentas ou processos manuais, nomeadamente a gestão de sócios, atletas, equipas, patrocinadores, eventos, bilheteira, pagamentos e documentos.

A aplicação resultante é composta por um \textit{backend} desenvolvido em Kotlin com Spring Boot e por um \textit{frontend} desenvolvido em React com TypeScript e Vite. Esta separação entre cliente e servidor permitiu definir responsabilidades claras: o \textit{backend} concentra a lógica de negócio, persistência, autenticação, autorização e integrações externas, enquanto o \textit{frontend} disponibiliza uma interface web para os diferentes perfis de utilização. A utilização de PostgreSQL como base de dados relacional, JDBI para acesso a dados, Stripe para pagamentos, Cloudflare R2 para armazenamento de ficheiros e Docker para execução local contribuiu para uma solução modular, extensível e próxima de um ambiente real de produção.

Um dos aspetos mais relevantes do projeto foi a organização modular do \textit{backend}. A separação entre os módulos de domínio, serviços, repositórios, camada HTTP e arranque da aplicação permitiu reduzir o acoplamento entre componentes e facilitar a testabilidade. Esta decisão revelou-se especialmente importante devido à dimensão funcional da aplicação, que envolve várias áreas de negócio relacionadas entre si. A utilização de interfaces para os repositórios e de um \texttt{TransactionManager} permitiu ainda isolar a lógica de negócio da tecnologia concreta de persistência.

No \textit{frontend}, a organização por domínio funcional permitiu manter próximos os componentes, tipos, estilos, \textit{hooks} e chamadas à API associados a cada área da aplicação. Esta abordagem facilitou a evolução da interface e permitiu criar fluxos distintos para utilizadores públicos, utilizadores autenticados e utilizadores com permissões administrativas. A existência de rotas protegidas, contexto de autenticação e internacionalização contribuiu para uma experiência mais completa e adequada ao contexto do clube.

Relativamente aos objetivos definidos, considera-se que os requisitos principais foram cumpridos. A aplicação permite a criação e gestão de sócios, inscrição e aprovação de atletas, gestão de patrocinadores e contratos de patrocínio, criação de eventos, compra de bilhetes, integração com pagamentos online, gestão de ficheiros e administração das principais entidades do sistema. Para além disso, foram produzidos elementos de apoio importantes, como documentação da API em OpenAPI, ficheiros de configuração para Docker e testes automatizados sobre componentes relevantes da aplicação.

Durante o desenvolvimento surgiram também desafios significativos. A modelação do domínio exigiu cuidado, uma vez que várias entidades possuem relações entre si e estados próprios, como sócios, atletas, patrocínios, cobranças e pagamentos. A integração com a Stripe exigiu uma atenção particular ao fluxo de pagamento, à validação de notificações externas e à atualização correta do estado interno da aplicação. Outro desafio foi garantir que a interface se mantinha utilizável apesar da quantidade de funcionalidades disponíveis, sobretudo na área administrativa.

Apesar de a aplicação já cobrir um conjunto alargado de funcionalidades, existem várias melhorias possíveis para trabalho futuro. Entre elas destacam-se o reforço da cobertura de testes, sobretudo em fluxos completos de integração, a melhoria da experiência de utilização em dispositivos móveis nas áreas administrativas, a adição de mecanismos avançados de pesquisa e filtragem, a criação de notificações automáticas para pagamentos ou aprovações pendentes, e a preparação de uma configuração de produção mais completa, incluindo monitorização, cópias de segurança e políticas de segurança adicionais.

Em suma, a JagozApp demonstra a viabilidade de uma plataforma integrada para apoiar a gestão de um clube desportivo local. O projeto permitiu aplicar conhecimentos de desenvolvimento web, arquitetura de software, bases de dados, autenticação, integração com serviços externos, testes e documentação. Mais do que uma aplicação isolada, o resultado final constitui uma base sólida que pode continuar a evoluir de acordo com as necessidades reais do clube.